\section{Formal Statements and Proofs}
\label{sec:appendix_proofs}

This appendix collects theorem-level formulations used in the main text.
The body of the manuscript keeps the same results in interpretation-first form.

\subsection{Strong-weak involution and asymptotic transfer}

\begin{theorem}[Strong-weak representative involution]
\label{thm:strong_weak_involution_app}
Assume one-channel scaling
\begin{equation}
  \chi(\beta,g,\vartheta)=g^2\chi_0(\beta,\vartheta),\qquad \chi_0(\beta,\vartheta)>0.
\end{equation}
Define
\begin{equation}
  g_\star(\beta,\vartheta):=\chi_0(\beta,\vartheta)^{-1/2},
  \qquad
  g^\vee:=\frac{g_\star^2}{g}=\frac{1}{\chi_0 g}.
\end{equation}
Then
\begin{align}
  \chi(\beta,g^\vee,\vartheta) &= \chi(\beta,g,\vartheta)^{-1}, \\
  (g^\vee)^\vee &= g, \\
  y(\beta,g^\vee,\vartheta) &= -y(\beta,g,\vartheta),\qquad y:=\log\chi.
\end{align}
\end{theorem}
\begin{proof}
Direct substitution gives
\begin{equation}
  \chi(\beta,g^\vee,\vartheta)
  =\chi_0\left(\frac{1}{\chi_0 g}\right)^2
  =\frac{1}{\chi_0 g^2}
  =\chi(\beta,g,\vartheta)^{-1}.
\end{equation}
Applying the map again,
\begin{equation}
  (g^\vee)^\vee=\frac{1}{\chi_0 g^\vee}=\frac{1}{\chi_0}\,(\chi_0 g)=g.
\end{equation}
Taking logarithms yields \(y\mapsto-y\).
\end{proof}

\begin{proposition}[Asymptotic exchange for admissible observables]
\label{prop:asymptotic_exchange_app}
Let \(Q\) satisfy Definition~\ref{def:admissible_class}, and let each building block
\(\phi_j\) admit weak and strong asymptotic expansions at \(\chi\to0\) and \(\chi\to\infty\).
Then the strong-\(\chi\) asymptotic expansion of \(Q\) at \((\beta,g,\vartheta)\) is obtained by
evaluating the weak-\(\chi\) expansion at \((\beta,g^\vee,\vartheta)\), and conversely.
\end{proposition}
\begin{proof}
By Theorem~\ref{thm:strong_weak_involution_app},
\(\chi(\beta,g^\vee,\vartheta)=\chi(\beta,g,\vartheta)^{-1}\).
Hence a weak expansion
\(\phi_j(\chi)\sim\sum_n a_{j,n}\chi^n\) at \(g^\vee\) becomes
\(\sum_n a_{j,n}\chi^{-n}\) at \(g\), i.e. a strong expansion.
The reverse direction is identical.
\end{proof}

\subsection{Orientation gauge and affine reflection}

\begin{lemma}[Affine effective-temperature reflection]
\label{lem:affine_reflection_app}
Let
\begin{equation}
  \beff=\beta-\delta\beta,
  \qquad
  \delta\beta:=\frac{2}{\omega_q}\operatorname{arctanh}(u),
\end{equation}
with \(|u|<1\).
Under orientation gauge \(u\mapsto-u\),
\begin{equation}
  \beff' = 2\beta-\beff.
\end{equation}
\end{lemma}
\begin{proof}
Because \(\operatorname{arctanh}\) is odd,
\begin{equation}
  \delta\beta'=
  \frac{2}{\omega_q}\operatorname{arctanh}(-u)
  =-\delta\beta.
\end{equation}
Therefore
\begin{equation}
  \beff' = \beta-\delta\beta' = \beta+\delta\beta = 2\beta-\beff.
\end{equation}
\end{proof}

\begin{lemma}[Even/odd observable split]
\label{lem:even_odd_split_app}
If \(O_{\mathrm{rad}}=F(\chi)\), then orientation gauge leaves it invariant.
If \(O_{\mathrm{odd}}=H(\chi,u)\) and \(H(\chi,-u)=-H(\chi,u)\), then
\(O_{\mathrm{odd}}\mapsto -O_{\mathrm{odd}}\).
\end{lemma}
\begin{proof}
Orientation gauge acts as \((\chi,u)\mapsto(\chi,-u)\).
Substituting into \(F\) and \(H\) gives the parity statements.
\end{proof}

\subsection{Composed action structure}

\begin{proposition}[Composed map and commuting radial action]
\label{prop:composed_action_app}
Let \(\mathcal{D}_{\mathrm{sw}}\) denote strong-weak reflection and
\(\mathcal{G}\) orientation gauge.
Both are involutions.
On radial sectors,
\begin{equation}
  (\mathcal{G}\circ\mathcal{D}_{\mathrm{sw}})\chi
  =(\mathcal{D}_{\mathrm{sw}}\circ\mathcal{G})\chi
  =\chi^{-1}.
\end{equation}
On oriented sectors they commute whenever strong-weak mapping preserves orientation parity.
\end{proposition}
\begin{proof}
Involution of \(\mathcal{D}_{\mathrm{sw}}\) follows from
Theorem~\ref{thm:strong_weak_involution_app}, and involution of \(\mathcal{G}\) is immediate from
sign flip.
Orientation acts trivially on \(\chi\), so radial commutation follows.
The parity-preserving condition is sufficient for commutation on oriented labels.
\end{proof}

\subsection{Normal-duality existence and involution}

\begin{definition}[Normal-dual map]
\label{def:normal_dual_app}
Fix \(\theta\), define \(y(\beta,g,\theta)=\log\chi(\beta,g,\theta)\), and
\(\mathcal{C}_\theta:=\{(\beta,g):y=0\}\).
Assume \(y\in C^2\), \(\nabla y\neq0\) on \(\mathcal{C}_\theta\), and a tubular neighborhood where each
point \(p\) has unique projection \(\pi(p)\in\mathcal{C}_\theta\).
Writing
\begin{equation}
  p=\pi(p)+\lambda\hat\nu(\pi(p)),
\end{equation}
with \(\hat\nu\) unit normal, define \(p^\perp\) on the same normal ray by
\begin{equation}
  y(p^\perp)=-y(p).
\end{equation}
\end{definition}

\begin{theorem}[Local existence and involution of normal duality]
\label{thm:normal_duality_app}
Under Definition~\ref{def:normal_dual_app}, there exists \(\delta>0\) such that for
\(|\lambda|<\delta\):
\begin{enumerate}
  \item \(p^\perp\) exists and is unique,
  \item the map is involutive, \((p^\perp)^\perp=p\),
  \item the segment connecting \(p\) and \(p^\perp\) is orthogonal to
  \(\mathcal{C}_\theta\) at \(\pi(p)\).
\end{enumerate}
\end{theorem}
\begin{proof}
For fixed \(c\in\mathcal{C}_\theta\), define
\(\Phi_c(\lambda):=y(c+\lambda\hat\nu(c))\).
Since \(\Phi_c(0)=0\) and
\(\Phi_c'(0)=\nabla y(c)\cdot\hat\nu(c)=\|\nabla y(c)\|>0\),
\(\Phi_c\) is strictly monotone on a sufficiently small interval.
Hence \(\Phi_c(\lambda^\perp)=-\Phi_c(\lambda)\) has a unique solution.
Applying the construction twice gives
\(\Phi_c(\lambda^{\perp\perp})=\Phi_c(\lambda)\), so
\(\lambda^{\perp\perp}=\lambda\) by injectivity.
Orthogonality is immediate because both points lie on one normal ray through \(c\).
\end{proof}

\subsection{Equation bookkeeping index}
Maintain one-to-one mapping with
\path{paper_dual_hmf/notes/equation_registry.md}.
